% rtex-literals.tex: readable PDF literal nodes, shared by capture and live compilation.
% Before input, define \rtex@litns as the Lua table providing literal(data).
\begingroup\catcode`\@=11\relax
% pgf's PDF literals, made from Lua so the display list can read their operators (see
% \rtex@litns.literal). \protected: like the primitive, never expanded inside an \edef.
\gdef\rtex@hookpgfliterals{%
  \ifdefined\pgfsys@invoke
    \protected\def\pgfsys@invoke##1{\directlua{\rtex@litns.literal("\luaescapestring{##1}")}}%
    \protected\def\pgfsys@color@unstacked##1{\directlua{\rtex@litns.literal("\luaescapestring{\csname\string\color@##1\endcsname}")}}%
  \fi}
% Also when the driver is read: pictures made in the preamble (the fadings library's fading
% boxes) are drawn by their literals too.
\AddToHook{file/pgfsys-luatex.def/after}{\rtex@hookpgfliterals}
\AddToHook{begindocument/before}{\rtex@hookpgfliterals}
% pict2e's PDF code (picture mode, diagbox), the same way: p2e-luatex.def writes
% \pdfextension literal{ q #1 Q }, a token list the display list cannot read.
% Only the driver's own definition is replaced (option `hide` gobbles the code).
\gdef\rtex@pIIe@driver#1{\pdfextension literal{ q #1 Q }}
\gdef\rtex@hookpicttwoe{%
  \ifx\pIIe@code\rtex@pIIe@driver
    \protected\def\pIIe@code##1{\directlua{\rtex@litns.literal("\luaescapestring{ q ##1 Q }")}}%
  \fi}
\AddToHook{package/pict2e/after}{\rtex@hookpicttwoe}
\endgroup
